\documentclass{fiche}
\metadonnees{14}{Une nuit de novembre}{Bloc 3 — La ville, le travail, la science}

\begin{document}
\entete

\objectifs{%
\textbf{Langue} : voix passive ; passif des verbes à deux compléments ; \emph{he is said
to}.\par
\textbf{Texte} : Mary Shelley, \emph{Frankenstein} (1818), chapitre \textsc{v}.\par
\textbf{Durée} : 1\,h\,30 — exercices 35 min, texte et questions 30 min, version 25 min.}

%% ===================================================================
\exercice[13 min]{Mettre au passif}

\rappel{\textbf{Rappel.} Le passif se forme avec \emph{be} + participe passé. Il met en tête
ce qui subit l'action ; l'agent n'apparaît après \emph{by} que s'il apporte une information.}

\consigne{Réécrivez au passif, en supprimant l'agent quand il n'apporte rien.}

\begin{anglais}
\begin{enumerate}[label=\arabic*.,itemsep=2.5mm,leftmargin=*]
  \item Somebody has opened the window.
    \lignes{1}
    \corr{The window has been opened.}
  \item They will finish the experiment tonight.
    \lignes{1}
    \corr{The experiment will be finished tonight.}
  \item Mary Shelley wrote \emph{Frankenstein} at nineteen.
    \lignes{1}
    \corr{\emph{Frankenstein} was written by Mary Shelley at nineteen. \emph{(l'agent est
    informatif : on le garde)}}
  \item People were watching the house all night.
    \lignes{1}
    \corr{The house was being watched all night.}
  \item Someone must clean these instruments.
    \lignes{1}
    \corr{These instruments must be cleaned.}
  \item Nobody had ever attempted such a thing.
    \lignes{1}
    \corr{Such a thing had never been attempted.}
  \item The rain was beating the panes.
    \lignes{1}
    \corr{The panes were being beaten by the rain.}
  \item They have not yet explained the accident.
    \lignes{1}
    \corr{The accident has not yet been explained.}
  \item One cannot undo this.
    \lignes{1}
    \corr{This cannot be undone.}
  \item The students are translating the text.
    \lignes{1}
    \corr{The text is being translated.}
\end{enumerate}
\end{anglais}

\note[Pourquoi le passif]{L'anglais y recourt beaucoup plus que le français, pour trois
raisons : l'agent est inconnu (1, 8), il est évident ou sans intérêt (5, 6, 10), ou bien on
veut garder le même sujet d'une phrase à l'autre. En version, on le rend le plus souvent par
\og on \fg{} + actif : \emph{the accident has not been explained} $\rightarrow$ \og on n'a
pas encore expliqué l'accident \fg.}

%% ===================================================================
\exercice[10 min]{Deux passifs, et le passif d'opinion}

\rappel{\textbf{Rappel.} Un verbe à deux compléments donne deux passifs, et l'anglais préfère
celui qui commence par la personne. \emph{It is said that he is\ldots} devient \emph{He is
said to be\ldots}}

\consigne{Complétez.}

\begin{anglais}
\begin{enumerate}[label=\arabic*.,itemsep=2.2mm,leftmargin=*]
  \item They gave him a prize. $\rightarrow$ He \trou{was given} a prize.
  \item Somebody told us the truth. $\rightarrow$ We \trou{were told} the truth.
  \item They offered her a place. $\rightarrow$ She \trou{was offered} a place.
  \item Nobody has shown me the letter. $\rightarrow$ I \trou{have not been shown} the letter.
  \item They asked me a strange question. $\rightarrow$ I \trou{was asked} a strange question.
  \item It is said that he works at night. $\rightarrow$ He \trou{is said to work} at night.
  \item It is believed that the creature escaped. $\rightarrow$ The creature \trou{is believed
    to have escaped}.
  \item It is thought that she is in Geneva. $\rightarrow$ She \trou{is thought to be} in
    Geneva.
  \item It is known that they were students. $\rightarrow$ They \trou{are known to have been}
    students.
  \item It is expected that the storm will last. $\rightarrow$ The storm \trou{is expected to
    last}.
\end{enumerate}
\end{anglais}

\note[La règle du décalage]{Quand le fait rapporté est antérieur à l'opinion, l'infinitif
passe au parfait : \emph{it is believed that the creature escaped} $\rightarrow$
\emph{the creature is believed to have escaped} (7, 9). Même mécanisme qu'avec les modaux
(fiche 8) : le verbe principal reste au présent, c'est l'infinitif qui recule.}

\note[Un tour sans équivalent]{\emph{He was given a prize} n'a pas de correspondant direct en
français : \og il fut donné un prix \fg{} est impossible. On passe par \og on lui a donné un
prix \fg. C'est le passif le plus déroutant en version, et le plus fréquent.}

%% ===================================================================
\exercice[12 min]{Le corps et l'horreur}
\consigne{a) Associez chaque mot à sa traduction.}

\begin{multicols}{2}
\begin{enumerate}[label=\arabic*.,itemsep=1.2mm,leftmargin=*]
  \item a limb \hfill \trou{g}
  \item a socket \hfill \trou{c}
  \item a complexion \hfill \trou{j}
  \item lustrous \hfill \trou{a}
  \item shrivelled \hfill \trou{e}
  \item a pane \hfill \trou{i}
  \item dreary \hfill \trou{b}
  \item to endure \hfill \trou{l}
  \item a shroud \hfill \trou{d}
  \item a wretch \hfill \trou{f}
  \item to pattern \hfill \trou{h}
  \item lassitude \hfill \trou{k}
\end{enumerate}
\columnbreak
\begin{enumerate}[label=\alph*.,itemsep=1.2mm,leftmargin=*]
  \item luisant, éclatant
  \item morne, lugubre
  \item une orbite (de l'œil)
  \item un linceul
  \item ratatiné, desséché
  \item un misérable, un malheureux
  \item un membre
  \item crépiter, tambouriner
  \item une vitre, un carreau
  \item le teint
  \item l'accablement, la lassitude
  \item supporter, endurer
\end{enumerate}
\end{multicols}

\note[Faux amis]{\emph{A complexion} est le teint, non la complexion ni la complexité.
\emph{To endure} est \og supporter \fg, non \og endurcir \fg. \emph{A demon} au sens ancien
est un être surnaturel, pas nécessairement le diable : d'où \emph{demoniacal},
\og démoniaque \fg.}

\consigne{b) Complétez avec un mot de la liste.}
\begin{anglais}
\begin{enumerate}[label=\arabic*.,itemsep=2mm,leftmargin=*]
  \item The rain \trou{pattered} dismally against the \trou{panes}.
  \item His eyes were set in \trou{sockets} of a dull white.
  \item Unable to \trou{endure} the sight of the \trou{wretch}, he rushed out of the room.
  \item A \trou{shroud} enveloped her form, and the grave-worms crawled in its folds.
\end{enumerate}
\end{anglais}

%% ===================================================================
\newpage
\section{Texte}

\noindent\small\textit{Victor Frankenstein, étudiant à Ingolstadt, travaille depuis deux ans
à donner la vie à un corps qu'il a assemblé lui-même. Il raconte la nuit où il y
parvient.}\normalsize

\begin{texte}
It was on a \gl[dreary]{dreary}{morne, lugubre} night of November that I
\gl[to behold, beheld]{beheld}{contempler, voir} the accomplishment of my
\gl[toils]{toils}{les peines, les labeurs}. With an anxiety that almost amounted to agony, I
collected the instruments of life around me, that I might infuse a spark of being into the
lifeless thing that lay at my feet. It was already one in the morning; the rain
\gl[to patter]{pattered}{crépiter, tambouriner} dismally against the
\gl[a pane]{panes}{une vitre}, and my candle was nearly burnt out, when, by the
\gl[a glimmer]{glimmer}{une lueur faible} of the half-extinguished light, I saw the dull
yellow eye of the creature open; it breathed hard, and a convulsive motion agitated its limbs.

How can I describe my emotions at this catastrophe, or how
\gl[to delineate]{delineate}{dépeindre, décrire} the \gl[a wretch]{wretch}{un misérable,
une créature pitoyable} whom with such infinite pains and care I had
\gl[to endeavour]{endeavoured}{s'efforcer de} to form? His limbs were in proportion, and I had
selected his features as beautiful. Beautiful! Great God! His yellow skin scarcely covered
the work of muscles and arteries beneath; his hair was of a
\gl[lustrous]{lustrous}{luisant, éclatant} black, and flowing; his teeth of a pearly
whiteness; but these \gl[luxuriances]{luxuriances}{ces richesses, ces splendeurs} only formed
a more horrid contrast with his watery eyes, that seemed almost of the same colour as the
\gl[dun-white]{dun-white}{d'un blanc terne} \gl[a socket]{sockets}{une orbite} in which they
were set, his \gl[shrivelled]{shrivelled}{ratatiné, desséché}
\gl[a complexion]{complexion}{le teint} and straight black lips.

The different accidents of life are not so changeable as the feelings of human nature. I had
worked hard for nearly two years, for the sole purpose of infusing life into an inanimate
body. For this I had deprived myself of rest and health. I had desired it with an
\gl[ardour]{ardour}{l'ardeur} that far exceeded moderation; but now that I had finished, the
beauty of the dream vanished, and breathless horror and disgust filled my heart.
\end{texte}
\source{Mary Shelley, \emph{Frankenstein; or, The Modern Prometheus}, Londres, Lackington,
1818, ch.~\textsc{v}.}

\partiel[15 min]{Questions}
\consigne{\en{Answer in English, in full sentences.}}

\begin{enumerate}[label=\arabic*.,itemsep=2mm,leftmargin=*]
  \item \en{What are the details of the setting, and what do they contribute?}
    \lignes{3}
    \corr{A dreary night of November, one in the morning, rain pattering dismally on the
    panes, a candle nearly burnt out. Everything is dying down at the moment when something
    is about to come to life, which is exactly the reversal the scene rests on.}
  \item \en{Which word does Victor use for the moment he has worked two years to reach, and
    why is it surprising?}
    \lignes{2}
    \corr{He calls it \emph{this catastrophe}. The word belongs to disaster, not to success:
    the narrator, writing afterwards, already knows what the triumph will cost.}
  \item \en{Victor says he had selected the features as beautiful. Why is the result
    horrible?}
    \lignes{3}
    \corr{Each part was chosen for its beauty --- proportioned limbs, lustrous flowing hair,
    pearly teeth --- but the yellow skin scarcely covers the muscles beneath, and the watery
    eyes are almost the colour of their sockets. It is the contrast between the beautiful
    parts and the dead ones that horrifies.}
  \item \en{What had the two years' work cost him?}
    \lignes{2}
    \corr{He had deprived himself of rest and health, and had desired his aim with an ardour
    that far exceeded moderation.}
  \item \en{``The beauty of the dream vanished.'' What exactly has changed?}
    \lignes{3}
    \corr{Nothing in the creature: it is Victor's feeling that has reversed the instant the
    work was finished. The sentence before says it: \emph{the different accidents of life are
    not so changeable as the feelings of human nature}.}
  \item \en{Whose side is the reader on at the end of this passage --- and why is that going
    to be a problem?}
    \lignes{3}
    \corr{We see everything through Victor's eyes and share his disgust. But the only thing
    the creature has done is open an eye and breathe, and the reader who notices that has
    already begun to read against the narrator.}
\end{enumerate}

\note[Les passifs du texte]{\emph{my candle was nearly burnt out}, \emph{his limbs were in
proportion}, \emph{his features\ldots I had selected}, \emph{the sockets in which they were
set}, \emph{the half-extinguished light} : le passif et le participe passé épithète servent
ici à effacer l'agent --- et l'agent, c'est Victor. La langue prend ses distances avant que le
personnage ne le fasse.}

%% ===================================================================
\newpage
\section{Version}
\consigne{Traduisez le passage en français. Il suit immédiatement le texte.}

\begin{version}
Unable to endure the aspect of the being I had created, I rushed out of the room and continued
a long time traversing my bed-chamber, unable to compose my mind to sleep. At length
\gl[lassitude]{lassitude}{l'accablement} succeeded to the tumult I had before endured, and I
threw myself on the bed in my clothes, endeavouring to seek a few moments of forgetfulness.
But it was in vain; I slept, indeed, but I was disturbed by the wildest dreams. I thought I
saw Elizabeth, in the bloom of health, walking in the streets of Ingolstadt. Delighted and
surprised, I embraced her, but as I imprinted the first kiss on her lips, they became
\gl[livid]{livid}{blême, plombé} with the \gl[a hue]{hue}{une teinte} of death; her features
appeared to change, and I thought that I held the corpse of my dead mother in my arms; a
\gl[a shroud]{shroud}{un linceul} enveloped her form, and I saw the grave-worms crawling in
the folds of the \gl[flannel]{flannel}{la flanelle}. I started from my sleep with horror; a
cold dew covered my forehead, my teeth chattered, and every limb became convulsed; when, by
the dim and yellow light of the moon, as it forced its way through the window shutters, I
beheld the wretch --- the miserable monster whom I had created. He held up the curtain of the
bed; and his eyes, if eyes they may be called, were fixed on me. His jaws opened, and he
muttered some inarticulate sounds, while a grin \gl[to wrinkle]{wrinkled}{plisser} his cheeks.
He might have spoken, but I did not hear; one hand was stretched out, seemingly to
\gl[to detain]{detain}{retenir} me, but I escaped and rushed downstairs.
\end{version}
\source{\emph{Ibid.}}

\lignes{22}

\corr{\textbf{Proposition de traduction.} Incapable de supporter la vue de l'être que j'avais
créé, je me précipitai hors de la pièce et arpentai longtemps ma chambre, sans pouvoir
calmer mon esprit et trouver le sommeil. À la fin, l'accablement succéda au tumulte que
j'avais enduré jusque-là, et je me jetai tout habillé sur le lit, m'efforçant de chercher
quelques instants d'oubli. Mais ce fut en vain : je dormis, certes, mais des rêves égarés
vinrent me troubler. Je crus voir Elizabeth, dans tout l'éclat de la santé, qui marchait dans
les rues d'Ingolstadt. Ravi et surpris, je l'embrassai ; mais comme je déposais le premier
baiser sur ses lèvres, elles se plombèrent de la teinte de la mort ; ses traits parurent
changer, et je crus tenir dans mes bras le cadavre de ma mère morte ; un linceul enveloppait
son corps, et je vis les vers du tombeau ramper dans les plis de la flanelle. Je m'éveillai
en sursaut, saisi d'horreur ; une sueur froide couvrait mon front, mes dents claquaient et
tous mes membres étaient pris de convulsions ; quand, à la lueur blafarde et jaune de la lune,
qui se frayait un passage à travers les volets, j'aperçus le misérable --- le monstre
lamentable que j'avais créé. Il soulevait le rideau du lit, et ses yeux, si l'on peut appeler
cela des yeux, étaient fixés sur moi. Ses mâchoires s'ouvrirent et il marmonna quelques sons
inarticulés, tandis qu'un rictus plissait ses joues. Peut-être parla-t-il, mais je n'entendis
pas ; une main se tendait, comme pour me retenir, mais je m'échappai et me précipitai en bas
de l'escalier.}

\note[Choix de traduction 1]{\emph{the being I had created} : relative sans pronom (fiche 13).
Trois autres dans le passage : \emph{the tumult I had before endured}, \emph{the monster whom
I had created}. Les repérer d'abord, traduire ensuite.}

\note[Choix de traduction 2]{\emph{I was disturbed by the wildest dreams} : passif avec agent.
Le français peut le garder, mais l'actif est plus naturel ici (\og des rêves égarés vinrent me
troubler \fg). Choisir selon le rythme, pas selon la grammaire.}

\note[Choix de traduction 3]{\emph{if eyes they may be called} : inversion emphatique héritée
du style biblique. Le français rétablit l'ordre et ajoute \og on \fg{} : \og si l'on peut
appeler cela des yeux \fg.}

\note[Choix de traduction 4]{\emph{He might have spoken, but I did not hear} : \emph{might
have} + participe, ici hypothèse sur le passé, non regret : \og peut-être parla-t-il \fg. Le
contexte tranche, comme en fiche 9.}

\note[Choix de traduction 5]{\emph{a grin} n'est pas un sourire mais un rictus, une grimace
qui découvre les dents. Faux ami dangereux, d'autant que le mot revient dans la bouche de
Victor chaque fois qu'il décrit la créature.}

%% ===================================================================
\partiel{Lexique à mémoriser}
\begin{itemize}[label=--,itemsep=0.8mm,leftmargin=*]\small
  \item \textbf{dreary} — morne, lugubre
  \item \textbf{a spark} — une étincelle
  \item \textbf{lifeless} — sans vie ; \textit{suffixe privatif} -less \textit{(fiche 2)}
  \item \textbf{a pane} — une vitre ; a window-pane
  \item \textbf{a glimmer} — une lueur faible ; \textit{verbe} to glimmer
  \item \textbf{a limb} — un membre \textit{(fiche 8)}
  \item \textbf{a feature} — un trait \textit{(fiche 11)}
  \item \textbf{a complexion} — le teint
  \item \textbf{a socket} — une orbite ; \textit{et} une prise de courant
  \item \textbf{a wretch} — un misérable ; wretched \textit{(fiche 7)}
  \item \textbf{to endure} — supporter
  \item \textbf{to deprive of} — priver de ; \textit{nom} deprivation
  \item \textbf{a purpose} — un but \textit{(fiche 7)} ; on purpose \textit{« exprès »}
  \item \textbf{to vanish} — s'évanouir, disparaître
  \item \textbf{a corpse} — un cadavre ; \textit{le} p \textit{se prononce}
  \item \textbf{a shroud} — un linceul
  \item \textbf{to crawl} — ramper
  \item \textbf{a grin} — un rictus, un large sourire forcé
  \item \textbf{to mutter} — marmonner
  \item \textbf{to stretch out} — tendre, étirer
\end{itemize}

\end{document}
